%% ============================================================
%%  Chapter 1 — Introduction
%% ============================================================
\chapter{Introduction}
\label{ch:introduction}

\section{Motivation and Context}

The modern networked enterprise operates in an environment of persistent,
structured adversarial pressure.
The global cost of cybercrime exceeded USD\,8 trillion in 2023 and is
projected to reach USD\,10.5 trillion annually by 2025~\citep{cisa2023threat}.
These figures reflect not merely the direct cost of successful intrusions ---
data theft, ransomware payments, system downtime --- but also the growing
investment required to maintain defensive infrastructure against an ever-more-
sophisticated threat landscape.

Contemporary adversaries operate with strategic intent.
Advanced Persistent Threat (APT) groups sponsored by nation-states, organised
criminal syndicates, and increasingly automated malware ecosystems conduct
multi-stage campaigns that unfold over days, weeks, or months.
These campaigns follow recognisable patterns: initial reconnaissance of the
target network, followed by weapon development, delivery of malicious payloads,
exploitation of discovered vulnerabilities, establishment of persistent
backdoors, command-and-control communication, and finally the execution of
the mission objective.
The Lockheed Martin Cyber Kill Chain~\citep{hutchins2011intelligence} formalises
this progression into seven sequential phases and has become the standard
reference model for structuring both threat intelligence and defensive
countermeasures.

\subsection{The Limitations of Static Defences}

Traditional network defences are designed around the principle of known-threat
identification.
Signature-based Intrusion Detection Systems (IDS) such as Snort and Suricata
maintain databases of known malicious patterns and alert when network traffic
matches a signature.
This approach is effective against catalogued threats but fundamentally blind to
novel or zero-day exploits~\citep{tavallaee2009detailed}.
When an adversary deploys a new technique not yet in the signature database,
the IDS produces no alert.

Fixed-threshold anomaly detectors address this limitation by flagging deviations
from a statistical baseline rather than matching known patterns.
However, the baseline calibration problem is severe in practice: thresholds
that are sensitive enough to detect subtle intrusions tend to generate
unmanageable volumes of false alarms, overwhelming security operations centre
(SOC) analysts.
Thresholds set conservatively to reduce false alarms miss the stealthy,
low-and-slow campaigns favoured by sophisticated adversaries.

Rule-based honeypot management systems face a related challenge.
A honeypot administrator must specify in advance how the system should respond
to each observed traffic pattern: what level of engagement warrants logging
versus blocking versus alerting.
These rules are effective when the anticipated threat matches the rule designer's
expectations, but adversaries adapt.
An attacker who discovers that certain traffic patterns trigger immediate
blocking may modify their approach to stay below the detection threshold
while still achieving their objectives --- and, distinctly, an attacker who
is \emph{identified once} and simply waits out a fixed observation window can
present as a first-time visitor on every subsequent visit, since a purely
rule-based system with no cross-session memory has no way to treat a known
repeat offender differently from a stranger.

\subsection{Honeypots as Adaptive Decoys}

Honeypots offer a qualitatively different defensive capability.
A decoy system placed within or at the perimeter of a network attracts
attacker attention by appearing to be a valuable, vulnerable target.
Any interaction with the honeypot is inherently suspicious: legitimate users
have no reason to access a system that serves no production purpose.
This property eliminates the false positive problem for the detection
\emph{event} itself --- but the question of \emph{how to respond} to a
detected interaction remains.

The response decision is non-trivial for several reasons.
First, not all honeypot interactions are equally dangerous.
A port scan that visits a honeypot is suspicious but low-severity; a
successful exploit that installs a backdoor is critical.
The response should be proportional: aggressive blocking of low-severity
interactions wastes the intelligence value of honeypot engagement and
may alert the attacker that they have been detected.
Second, the optimal response depends on the attacker's position in the kill
chain: intelligence gathering at RECONNAISSANCE stage is best served by
quiet logging, while command-and-control activity at Stage 5 warrants
immediate escalation.
Third, the response should account for the attacker's likely \emph{next}
move: if the attacker is likely to escalate to a more dangerous stage,
the defender should pre-empt that escalation rather than waiting for it to
occur.
Fourth, and less obviously, the response should account for the source's
\emph{history}: an IP address that has previously conducted a sustained
attack campaign against the honeypot represents a materially different risk
than an IP address making its first-ever contact, even if the two IPs'
current traffic looks identical.

\subsection{Reinforcement Learning for Adaptive Defence}

Reinforcement learning (RL) provides a principled framework for learning
adaptive responses from experience~\citep{sutton2018reinforcement}.
An RL agent interacts with a simulated environment, observes the consequences
of its actions, and adjusts its policy to maximise cumulative reward.
Applied to honeypot management, an RL agent could potentially discover
response strategies that generalise across attacker behaviours and
optimally balance intelligence gathering against containment.

Deep Q-Networks (DQN)~\citep{mnih2015humanlevel} extend Q-learning to
high-dimensional state spaces by using a neural network to approximate
the action-value function.
DQN demonstrated human-level performance on Atari games and has subsequently
been applied to a range of sequential decision problems including robotic
control, traffic management, and network defence~\citep{elderman2017adversarial}.
However, DQN presents two significant challenges for operational cybersecurity:

\begin{enumerate}
  \item \textbf{Opacity}: The Q-values produced by a multi-layer perceptron
        cannot be directly audited by a human analyst.
        When a DQN agent chooses to BLOCK rather than LOG a specific
        traffic pattern, there is no traceable explanation.
        Post-hoc explanation methods such as SHAP~\citep{lundberg2017unified}
        provide approximations, but these are not the policy itself.
  \item \textbf{Data dependency}: DQN requires a substantial corpus of
        interaction data before converging to a useful policy.
        In a real honeypot deployment, collecting sufficient labelled
        attack data to train a DQN without the simulation-to-reality
        gap is itself a major challenge.
\end{enumerate}

This thesis implements and trains a DQN baseline against exactly this
system (Section~\ref{sec:dqn_baseline}, Section~4.5) and finds, empirically
rather than merely by citing the general literature, that both challenges
materialise concretely: the trained agent's decisions are unauditable by
construction, and its training data dependency produces a policy that
specialises to a narrow slice of the training distribution rather than
learning genuine threat discrimination (Section~5.2).
This finding is revisited directly in Section~5.3 when considering whether
a learned policy would be the appropriate mechanism for the dynamic
behavioural adjustments introduced later in this thesis
(Sections~\ref{sec:escalation_tracking}--\ref{sec:dynamic_response}).

\subsection{The Case for Interpretable Policies}

In high-stakes operational domains, interpretability is not optional.
Security operations centres are subject to regulatory requirements for
incident reporting; the decision to block an IP address or raise a
major alert must be defensible to auditors, management, and sometimes
courts.
A policy that cannot be explained is a policy that cannot be trusted with
critical decisions.

\citet{rudin2019stop} argues forcefully that in such domains, interpretable
models should be preferred over black-box models: not only are they more
auditable, they are often comparably accurate when the problem structure
permits a good interpretable solution.
The cybersecurity kill chain provides exactly this structure: the seven-stage
model defines a clear ordinal risk progression, and the Markov chain provides
exact, computable escalation probabilities.
These quantities are sufficient to define a transparent policy that captures
the key decision dimensions without requiring opaque function approximation.
This thesis extends that argument beyond the \emph{initial} decision policy
to the question of \emph{dynamic adjustment}:
Sections~\ref{sec:escalation_tracking}--\ref{sec:dynamic_response} show that
a honeypot's behaviour can be made to respond to sustained, observed traffic
patterns --- remembering a repeat offender, avoiding alert fatigue --- using
the same auditable-by-construction design philosophy, rather than reaching
for a learned adjustment mechanism by default.

\section{Problem Statement}

This thesis addresses four interrelated challenges in adaptive honeypot
management.

\paragraph{Challenge 1: Realistic attacker simulation.}
Any evaluation of a honeypot policy requires a realistic adversary.
Replaying captured attack logs does not support counterfactual evaluation:
if the defender changes its responses, a static replay cannot model how the
attacker would adapt.
A synthetic attacker must produce plausible network traffic, follow a coherent
escalation logic, and support multiple qualitatively distinct intent profiles
that exercise the defender across a range of threat scenarios.
The attacker must be parameterisable (to allow controlled experiments) and
grounded in real-world data (to maintain ecological validity).
Realism extends beyond the attacker: a defender evaluated only against
traffic where every session behaves identically to every other session of
the same type, and where benign traffic is a single homogeneous distribution,
risks learning --- or, for an interpretable policy, being \emph{validated
against} --- a narrower operating picture than a real deployment would
present.

\paragraph{Challenge 2: Intent-aware adaptive defence.}
A static rule base is insufficient for adaptive adversaries.
An effective honeypot policy must map observable signals --- current attack
type, kill chain stage, escalation rate --- to appropriate responses in a
way that generalises across different attacker behaviours and profiles.
The policy must balance \emph{sensitivity} (responding to genuine threats)
against \emph{specificity} (avoiding false alarms on borderline traffic)
across a range of threat-level distributions.

\paragraph{Challenge 3: Interpretability without sacrificing effectiveness.}
RL policies trained with DQN or similar algorithms achieve strong empirical
performance but offer no insight into their decision-making.
Operational deployment requires policies that can be inspected, audited, and
understood by human practitioners without relying on post-hoc approximations.
The challenge is to design a policy that is simultaneously interpretable,
computationally tractable, and effective across multiple attacker profiles.

\paragraph{Challenge 4: Dynamic behavioural adjustment without sacrificing
auditability.}
A policy fixed at design time cannot account for information that only
becomes available at runtime: whether a specific source has attacked before,
or whether a given override rule is firing so frequently that it has stopped
conveying useful signal to a downstream analyst.
Addressing this naively --- for example, by reaching directly for a learned
policy, as Chapter~\ref{ch:background} and Chapter~\ref{ch:discussion}
discuss --- risks reintroducing the opacity and data-dependency problems
that motivated the interpretable design in the first place (Challenge 3).
The challenge is therefore to introduce genuine, observation-driven dynamic
adjustment while preserving the property that every resulting decision
remains traceable to an explicit, auditable rule.

\section{Research Objectives}

This thesis addresses the following primary research objectives:

\begin{enumerate}
  \item \textbf{Attacker modelling.}
        Design a Markov-chain attacker grounded in the Lockheed Martin
        Cyber Kill Chain, with per-intent transition modifiers and
        UNSW-NB15-inspired parametric feature distributions, capable of
        producing qualitatively distinct attack campaigns from a shared
        base model.

  \item \textbf{SEDM design.}
        Develop the Stage-Escalation Decision Matrix (SEDM), a
        transparent, deterministic honeypot policy that maps kill chain
        stage and Markov-chain-derived escalation risk to a ranked
        response action, with override rules for high-impact attack types,
        elevated attack frequency, and --- as developed later in this
        thesis --- persistent cross-session reputation.

  \item \textbf{Threat quantification.}
        Define a composite threat-level metric that integrates multiple
        attack signals (type severity, kill chain stage, escalation rate,
        cumulative pressure) into a single normalised index that drives
        the reward function and the SEDM composite risk score.

  \item \textbf{DQN baseline.}
        Implement a DQN-based defender as a learned baseline, characterise
        its training dynamics, and compare its performance properties
        with the SEDM in terms of detection rate, false-positive rate,
        and interpretability.

  \item \textbf{Empirical evaluation.}
        Evaluate both policies using detection rate, false-positive rate,
        episode reward, and action distribution across all four attacker
        intents, and characterise the interpretability--performance
        trade-off in the context of operational cybersecurity requirements.

  \item \textbf{Traffic and session realism.}
        Extend the attacker and event-generation subsystems so that
        synthetic sessions have a persistent, session-level character and
        benign traffic is drawn from multiple distinct behavioural
        personas, and quantify the effect of this realism on classifier
        accuracy and downstream SEDM metrics.

  \item \textbf{Escalation-signal design.}
        Design and empirically compare a severity-weighted alternative to
        the sliding-window escalation-rate signal that drives the SEDM's
        R3 override and composite risk score.

  \item \textbf{Auditable dynamic response.}
        Design, implement, and empirically evaluate two non-learning
        mechanisms --- a cross-session reputation override and a bounded
        threshold controller --- that allow the SEDM's behaviour to
        respond to sustained, observed patterns at runtime while remaining
        fully traceable, and explicitly re-examine whether a
        learning-based mechanism would be the more practical choice for
        this objective, using this project's own DQN results as direct
        evidence.
\end{enumerate}

\section{Principal Contributions}

The principal contributions of this work are as follows:

\paragraph{HoneyIQ simulation framework.}
A fully self-contained attacker--defender simulation built on the Gymnasium RL
interface~\citep{towers2024gymnasium}.
The framework implements a 24-dimensional state space combining one-hot encoded
attack type, kill chain stage, and attacker intent with continuous threat level,
attack count, and escalation rate.
A five-action honeypot policy interface (ALLOW, LOG, TROLL, BLOCK, ALERT) is
connected to a domain-informed reward function with three layers of domain
knowledge.
A comprehensive metrics and visualisation layer enables reproducible evaluation
across all four attacker intents.
All results can be reproduced without access to live network data.

\paragraph{Intent-aware Markov attacker with session-coherent traffic
realism.}
A coupled pair of Discrete-Time Markov Chains governing attack-type and
kill-chain-stage transitions.
Four intent profiles (Stealthy, Aggressive, Targeted, Opportunistic) are
encoded as element-wise transition modifiers that produce qualitatively
different attack campaigns from a shared base transition structure.
Network features are sampled from per-attack parametric distributions
calibrated to reproduce UNSW-NB15 attack signatures.
Building on this, each simulated session now draws a persistent intensity
and (for benign traffic) a behavioural persona once per episode, so a
session's volume-shaped features read as one coherent source rather than
independently-resampled noise at every step, and benign traffic is
represented by three distinct personas rather than a single distribution
(\S\ref{sec:realism}).

\paragraph{Stage-Escalation Decision Matrix (SEDM) with a fourth override
rule.}
An interpretable, deterministic, zero-training-overhead policy that:
(1) computes escalation risk analytically from the intent-specific Markov
transition matrix; (2) classifies the risk into three bands;
(3) performs a $7 \times 3$ matrix lookup; and
(4) applies four prioritised override rules, extending the original three
with a cross-session reputation override (R4, \S\ref{sec:dynamic_response})
evaluated ahead of the unconditional-allow rule for normal traffic.
The SEDM is fully traceable from observable inputs to final action,
printable on a single page, and modifiable by hand without retraining
any model.

\paragraph{Composite threat-level metric.}
A weighted combination of attack severity (45\%), kill chain stage weight
(35\%), escalation rate (15\%), and cumulative attack count (5\%) partitioned
into five threat bands, driving both the reward function and the SEDM
composite risk score.
The weight assignment is grounded in domain reasoning about the relative
informational value of each signal.

\paragraph{Severity-weighted escalation tracking.}
An exponential-moving-average alternative to the original hard
sliding-window escalation-rate signal, using the same per-attack-type
severity weights already defined for the composite threat level, computed
alongside the original signal so the two can be compared directly rather
than substituted blindly (\S\ref{sec:escalation_tracking}, Section~4.9).

\paragraph{Auditable dynamic response mechanisms.}
A cross-session \texttt{ReputationTracker} that maintains a
half-life-decayed offense score per source address, independent of
individual session lifetimes, feeding the new R4 override rule; and a
bounded deadband controller (\texttt{AdaptiveThresholds}) that keeps the R3
override's trigger frequency near an operator-specified target,
deliberately scoped and documented as an alert-fatigue safety valve rather
than a correctness mechanism, since no ground-truth feedback signal exists
in this system to justify a stronger claim (\S\ref{sec:dynamic_response},
Sections~4.11--4.12).

\paragraph{A direct, evidence-grounded answer to whether a learning-based
mechanism should have been used instead.}
Rather than treating the choice between rule-based and learned dynamic
adjustment as an open question answered by appeal to the general
literature, this thesis uses its own DQN baseline's documented failure mode
(Section~5.2) and its own evaluation harness's defect history to argue,
with direct project-specific evidence, that the auditable mechanisms above
are the more practical choice at this time, and identifies precisely what
would need to exist before that conclusion should be revisited
(Section~5.3).

\paragraph{Empirical evaluation across four intent profiles, extended to the
new mechanisms.}
Across 30 evaluation episodes per attacker intent (120 total, 24\,000 steps),
the SEDM achieves detection rates of 99.93\% (Stealthy), 100.0\% (Aggressive),
99.97\% (Targeted), and 99.97\% (Opportunistic) under the original traffic
generator, with a measured false positive rate of 0.00\% for all four
intents under classifier-driven decisions.
Chapter~\ref{ch:results} additionally reports: the effect of the realism
extensions on classifier accuracy and the classifier-driven false-positive
rate; the R3 trigger-rate difference between window and EMA escalation
tracking; the offense count at which cross-session reputation first
overrides the unconditional-allow rule; and the bounded threshold
controller's behaviour under both saturated and mixed-traffic conditions.

\section{Scope and Limitations}

The scope of this thesis is intentionally bounded to enable rigorous evaluation:

\begin{itemize}
  \item The environment is a simulation; results are not validated against
        live network traffic or real honeypot deployments.
  \item Feature distributions are parametric approximations inspired by
        UNSW-NB15 rather than fits to the actual dataset records.
  \item The attacker's Markov chain evolves independently of the defender's
        actions; the defender cannot deter or redirect the attacker.
  \item The primary evaluation reports both an oracle mode (the SEDM decides
        from ground-truth attack labels, an upper bound) and a
        classifier-driven mode (the SEDM decides from the Random Forest
        classifier's predicted attack type, scored against ground truth);
        the classifier itself is trained and evaluated only on synthetic
        feature distributions.
  \item Network topology is abstracted; all interactions are single-session
        point contacts without spatial or topological structure.
  \item The cross-session reputation mechanism (\S\ref{sec:dynamic_response})
        has no persistent-identity concept in the synthetic Gymnasium
        evaluation environment (only the live OpenCanary-integration
        pipeline tracks source addresses across sessions), so its
        evaluation (Section~4.11) uses a dedicated simulated-visit protocol
        rather than the primary 30-episode evaluation harness.
  \item The bounded threshold controller (\S\ref{sec:dynamic_response}) is
        evaluated against synthetic escalation-rate traces rather than
        against a live analyst-labelled alert stream, since no such
        labelling pathway exists in this system (Section~5.3).
\end{itemize}

These limitations are acknowledged in Chapter~\ref{ch:discussion} and
form the basis for future work directions in Section~\ref{sec:future}.

\section{Thesis Outline}

The remainder of this thesis is organised as follows.

\noindent\textbf{Chapter~\ref{ch:background}} surveys the theoretical and
empirical foundations of HoneyIQ: reinforcement learning and Markov Decision
Processes, Deep Q-Networks, Discrete-Time Markov Chains, the Lockheed Martin
Cyber Kill Chain, honeypot technology, the UNSW-NB15 dataset, Random Forests,
interpretable machine learning, and --- new in this version --- reputation-based
access control and bounded control as design references for the dynamic
response mechanisms introduced in Chapter~\ref{ch:methodology}.
Related work in RL-based intrusion response, attacker simulation, and
honeypot optimisation is reviewed.

\noindent\textbf{Chapter~\ref{ch:methodology}} describes the complete system
design: the attacker module (Markov chains, intent profiles, feature
simulation, and the session-coherent intensity/persona extension), the
composite threat-level metric and reward function, the SEDM policy design
and its now four-step override sequence, the DQN baseline architecture
and training protocol, the Gymnasium environment, the Random Forest
classifier, the severity-weighted escalation-tracking alternative, the
cross-session reputation and bounded threshold-adaptation mechanisms, and
the evaluation infrastructure.

\noindent\textbf{Chapter~\ref{ch:results}} presents quantitative results
from eight sets of experiments: (1) SEDM cross-intent evaluation,
(2) per-intent action distribution analysis, (3) composite risk and
escalation risk characterisation, (4) DQN training dynamics,
(5) parameter-selection and robustness analysis, (6) the effect of
traffic-realism extensions on classifier accuracy and classifier-driven
false positives, (7) window vs.\ EMA escalation-tracking comparison, and
(8) cross-session reputation and bounded threshold-adaptation behaviour.
All numerical results are derived from the experiment logs in
\texttt{results/} and \texttt{logs/}.

\noindent\textbf{Chapter~\ref{ch:discussion}} interprets the experimental
findings, analyses the interpretability--performance trade-off, examines the
structural reasons for the SEDM's cross-intent consistency, discusses
limitations of the simulation, reviews key design decisions, directly
re-examines whether the new dynamic-response mechanisms should instead have
been learning-based, and outlines future work directions.

\noindent\textbf{Chapter~\ref{ch:conclusion}} summarises the thesis
contributions and key findings.

\noindent Appendix~\ref{app:api} provides the complete public API reference
for each module, including the extensions introduced in this version.
Appendix~\ref{app:hyperparams} lists all configurable parameters used in
training and evaluation, with their default values and rationale, including
the new traffic-realism, escalation-tracking, and dynamic-response
parameters.
